\documentclass[12pt,a4paper]{article}
\usepackage{amsmath}
\usepackage{amssymb}  % Provides \mathbb and other symbols
\usepackage{graphicx}
\usepackage{subcaption}
\usepackage{enumitem}
\usepackage{listings}
\usepackage{hyperref}
\usepackage{ulem}
\usepackage[table]{xcolor} % For coloured boxes and table cells
\usepackage{amsthm}
\usepackage{tcolorbox}
\usepackage{array}
\usepackage{mdframed}
\usepackage{cancel}
\usepackage[margin=2.5cm]{geometry}
\setlength{\parindent}{0pt}
\setlength{\parskip}{0pt}
\begin{document}
\begin{titlepage}
\centering
\vspace*{\fill}
{\Huge Algebra} \\[0.5cm]
{\Huge Home Work 6}\\[1.5cm]
{\Large Student: 26-712-224}\\[0.5cm]
{\Large \today}\par
\vspace{2cm}
{\normalsize "I wrote these solutions fully myself und would like to receive feedback on these"\par}
\vspace*{\fill}
\end{titlepage}
% ============================================================
% Question A
% ============================================================
\newpage
\noindent
\textbf{A)} Let $p$ be a prime number. We will give an easy proof of the fact that
 every group of order $p^2$ is abelian using that every nontrivial $p$-group has
 a nontrivial center. (Compare with Exercise B of Sheet 4.) Let $G$ be a
 $p$-group of order $p^2$ and assume that $G$ is non-abelian. Obtain a
 contradiction as follows.\\[0.6em]
(1) Show that $Z(G)$ is cyclic of order $p$.\\
(2) Fix $g\in G-Z(G)$. Show that $\operatorname{ord}(g)=p$.\\
(3) Show that $\langle g\rangle\cap Z(G)=\{1\}$.\\
(4) Show that $\langle g\rangle\cdot Z(G)=G$.\\
(5) Deduce that $G$ is abelian, and conclude.
% Write your solution to Question A here.
% ============================================================
% Question B
% ============================================================
\newpage
\noindent
\textbf{B)} Let $G$ be a finite group acting on a finite set $X$. For $g\in G$, we set
\begin{align*}
  X^g=\{x\in X\mid g\cdot x=x\}.
\end{align*}
This is the subset of elements of $X$ fixed by $g$. The goal of the exercise
is to prove the following formula, due to Burnside:
% S^g is retained exactly as printed in the supplied question sheet.
\begin{align*}
  \operatorname{card}(G\backslash X)
  =\sum_{g\in G}\frac{\operatorname{card}(S^g)}{\operatorname{card}(G)}.
  \tag{$\star$}
\end{align*}
\textit{Step 1.} We first consider the case where $X=G/H$, for a subgroup
$H<G$, endowed with the natural action given by
\begin{align*}
  G\times G/H&\longrightarrow G/H\\
  (g,xH)&\longmapsto gxH.
\end{align*}
To treat this case, we introduce the set
\begin{align*}
  E=\{(g,k)\in G\times G\mid k^{-1}gk\in H\}.
\end{align*}
We have two maps $p_1,p_2:E\to G$ given by $p_1(g,k)=g$ and $p_2(g,k)=k$.\\[0.6em]
(1) Show that $p_2$ is surjective. For $k\in G$, show that there is a
  bijection from $p_2^{-1}(k)$ to the subgroup $k\cdot H\cdot k^{-1}$.
  Deduce that
  \begin{align*}
    \operatorname{card}(E)=\operatorname{card}(H)\cdot\operatorname{card}(G).
  \end{align*}
  (\textit{Hint:} Use that $E$ is the disjoint union of the $p_2^{-1}(k)$'s.)\\[0.6em]
(2) Let $g\in G$. Show that there is a bijection from $p_1^{-1}(g)$ to the set
  \begin{align*}
    F_g&=\{k\in G\mid g\cdot(kH)=kH\}\\
       &=\{k\in G\mid kH\in(G/H)^g\}.
  \end{align*}
  Use that the $H$-cosets belonging to $(G/H)^g$ form a partition of the set
  $\{k\in G\mid kH\in(G/H)^g\}$ to obtain the equality
  \begin{align*}
    \operatorname{card}(F_g)
    =\operatorname{card}((G/H)^g)\cdot\operatorname{card}(H).
  \end{align*}
  Deduce that
  \begin{align*}
    \operatorname{card}(E)
    =\sum_{g\in G}\operatorname{card}((G/H)^g)\cdot\operatorname{card}(H).
  \end{align*}
(3) Use (1) and (2) to obtain the Burnside formula $(\star)$ for $X=G/H$
  endowed with the natural action of $G$.\\[0.6em]
\textit{Step 2.} We now prove the general case of the Burnside formula.
Let $O_1,\cdots,O_r$ be the orbits of the action of $G$ on $X$.
(Thus $G\backslash X=\{O_1,\cdots,O_r\}$.)\\[0.6em]
(4) Note that the group $G$ acts on every $O_i$. Show that $S^g$ is the
  disjoint union of the $(O_i)^g$'s. Deduce that it suffices to prove the
  Burnside formula for each one of the orbits $O_i$ endowed with its natural
  action of $G$.\\[0.6em]
(5) Let $s\in O_i$ and set $H=\operatorname{Stab}_G(s)$.
  Using the bijection $G/H\xrightarrow{\sim}O_i$, deduce a bijection between
  $(O_i)^g$ and $(G/H)^g$. Use Step 1 to conclude.
% Write your solution to Question B here.
\end{document}